\section{Mineurs de graphes}\label{sec:mineurs-de-graphes}

\begin{td-exo}[Ordres graphiques]\, % 1
    \begin{enumerate}
        \item Montrer que la relation de sous-graphe forme un ordre sur les graphes.
        \item On d\'efinit la relation suivante sur les graphes :
        \begin{equation*}
            G \ll G' \text{ si } n(G) \leq n(G') \quad\text{ et } m(G) \leq m(G').
        \end{equation*}
        La relation \(\ll\) forme-t-elle un ordre sur les graphes ?
    \end{enumerate}
\end{td-exo}

\begin{td-sol}[]\, % 1
    \begin{enumerate}
        \item On rappelle la d\'efinition d'un ordre :
        
        \begin{definition}
            Soit \(E\) un ensemble et \(\mathcal{R}\) une relation entre deux \'el\'ements de \(E\).
            On dit que \(\mathcal{R}\) est une \defemph{relation d'ordre} sur \(E\) si elle v\'erifie les propri\'et\'es suivantes :
            \begin{enumerate}[label=\roman*)]
                \item R\'eflexivit\'e : pour tout \(x \in E\), on a \(x\mathcal{R}x\).
                \item Antisym\'etrie : pour tous \(x,y\in E\), si \(x\mathcal{R}y\) et \(y\mathcal{R}x\) alors \(x = y\).
                \item Transitivit\'e : pour tous \(x,y,z\in E\), si \(x\mathcal{R}y\) et \(y\mathcal{R}z\) alors \(x\mathcal{R}z\).
            \end{enumerate}
        \end{definition}

        V\'erifions alors ces propri\'et\'es sur la relation de sous-graphe :
        \begin{enumerate}
            \item R\'eflexivit\'e : pour tout graphe \(G\) on a bien \(V(G) \subseteq V(G)\) et \(E(G) \subseteq E(G)\) donc \(G\) est bien sous-graphe de \(G\).

            \item Antisym\'etrie : soient \(G\) et \(H\) deux graphes tels que \(G\) est sous-graphe de \(H\) et \(H\) est sous-graphe de \(G\).
            On a :
            \begin{equation*}
                \begin{aligned}
                    \hphantom{\implies} & V(G) \subseteq V(H) \quad \text{ et } V(H) \subseteq V(G) \\
                    \implies & V(G) \subseteq V(H) \subseteq V(G) \\
                    \implies & V(G) = V(H).
                \end{aligned}
            \end{equation*}

            De m\^eme, on a :
            
            \begin{equation*}
                \begin{aligned}
                    \hphantom{\implies} & E(G) \subseteq E(H) \quad \text{ et } E(H) \subseteq E(G) \\
                    \implies & E(G) \subseteq E(H) \subseteq E(G) \\
                    \implies & E(G) = E(H).
                \end{aligned}
            \end{equation*}

            Ainsi, comme \(V(G) = V(H)\) et \(E(G) = E(H)\), on a bien \(G = H\).

            \item Transitivit\'e : soient \(G, H\) et \(I\) trois graphes tels que \(G\) est sous-graphe de \(H\) et \(H\) est sous-graphe de \(I\).
            Alors, on a :
            \begin{equation*}
                \begin{aligned}
                    \hphantom{\implies} & V(G) \subseteq V(H) \quad \text{ et } V(H) \subseteq V(I) \\
                    \implies & V(G) \subseteq V(H) \subseteq V(I) \\
                    \implies & V(G) \subseteq V(I).
                \end{aligned}
            \end{equation*}

            De m\^eme, on a :
            
            \begin{equation*}
                \begin{aligned}
                    \hphantom{\implies} & E(G) \subseteq E(H) \quad \text{ et } E(H) \subseteq E(I) \\
                    \implies & E(G) \subseteq E(H) \subseteq E(I) \\
                    \implies & E(G) \subseteq E(I).
                \end{aligned}
            \end{equation*}

            Ainsi, comme \(V(G) \subseteq V(I)\) et \(E(G) \subseteq E(I)\), on a bien \(G\) sous-graphe de \(I\).
        \end{enumerate}

        \item On pose \(G = ([a, b, c, d], [(a, b), (c, d)])\) et \(H = ([a, b, c, d], [(a, b), (b, c)])\) comme dans la figure ci-dessous :
        
        \vspace{2mm}
\ffigbox[\FBwidth]{%
\caption{\centering Illustration des graphes \(G\) et \(H\)}\label{fig:td1_ex1_fi1}
}{%
\tikzset{
    main node/.style={
    circle, 
    draw, 
    fill=blue!20, 
    inner sep=1pt, 
    font=\scriptsize, 
    minimum size=6mm, 
    text=black
    }
}%
\begin{subfigure}[b]{0.4\linewidth}
    \centering
    \begin{tikzpicture}[scale=2]
    \node[main node] (1) at (0,0) {\(a\)};
    \node[main node] (2) at (1,0) {\(b\)};
    \node[main node] (3) at (0,-1) {\(c\)};
    \node[main node] (4) at (1,-1) {\(d\)};
    \draw (1) -- (2);
    \draw (3) -- (4);
    \end{tikzpicture}
    \caption{Graphe \(G\)}\label{fig:td1_ex1_fi1_sf1} 
\end{subfigure}%
\hfill
\begin{subfigure}[b]{0.4\linewidth}
    \centering
    \begin{tikzpicture}[scale=2]
    \node[main node] (1) at (0,0) {\(a\)};
    \node[main node] (2) at (1,0) {\(b\)};
    \node[main node] (3) at (0,-1) {\(c\)};
    \node[main node] (4) at (1,-1) {\(d\)};
    \draw (1) -- (2);
    \draw (2) -- (3);
    \end{tikzpicture}
    \caption{Graphe \(H\)}\label{fig:td1_ex1_fi1_sf2}
\end{subfigure}%
}

        On a \(n(G) = n(H)\) et \(m(G) = m(H)\) donc \(G \ll H\) et \(H \ll G\).
        Pourtant, \(G \neq H\) donc la propri\'et\'e d'antisym\'etrie n'est pas respect\'ee et \(\ll\) ne d\'efinit pas un ordre sur les graphes.
    \end{enumerate}
\end{td-sol}

\begin{td-exo}[Sous-graphe induit dense]\, % 2
    Montrer que la borne donn\'ee par le th\'eor\`eme du sous-graphe induit dense est optimale.
\end{td-exo}

\begin{td-sol}\,\\
    On rappelle le th\'eor\`eme :
    \begin{theorem}[Sous-graphe induit dense]\,\\
        Si le degr\'e moyen d'un graphe \(G\) est \(\geq 2k\), alors \(G\) contient un sous-graphe induit \(F\) de degr\'e minimum \(\geq k + 1\).
    \end{theorem}

    Soit \(G\) un cycle de longueur \(n\).
    Alors, on a \(\ol d(G) = 2k = 2 \times 1\).
    On peut bien trouver un sous-graphe induit \(F\) de degr\'e minimum \(\geq k + 1 = 2\) (par exemple, le cycle lui-m\^eme) mais il n'existe pas de sous-graphe induit \(F\) de degr\'e minimum \(\geq k + 2 = 3\), car tous les sommets du cycle ont un degr\'e de \(2\).

    La borne est donc optimale.

    On peut aussi construire un exemple plus g\'en\'eral.
    On construit \(G\) de la mani\`ere suivante :
    \begin{itemize}
        \item On prend un graphe complet \(K_{k + 1}\) de degr\'e moyen \(k\).
        \item On ajoute au graphe un stable de taille \(M\) et on cr\'ee le biparti complet entre le stable et le graphe complet.
    \end{itemize}
    On obtient alors un graphe \(G\) de degr\'e total
    \begin{equation*}
        \left(k+1\right)\left(k+M\right) + M\left(k+1\right) = \left(k+1\right)\left(k+2M\right)
    \end{equation*}
    et de degr\'e moyen
    \begin{equation*}
        \frac{\left(k+1\right)\left(k+2M\right)}{k+1+M}.
    \end{equation*}
    En d\'eveloppant, on trouve que pour \(M \geq \frac{k^2 + k}{2}\), on a bien \(\ol d(G) \geq 2k\).
    Comme le graphe complet \(K_{k + 1}\) est un sous-graphe induit de \(G\), on a bien un sous-graphe induit de degr\'e minimum \(\geq k + 1\), mais il n'existe pas de sous-graphe induit de degr\'e minimum \(\geq k + 2\), car tous les sommets du stable ont un degr\'e de \(k + 1\).
\end{td-sol}

\begin{td-exo}[Gros sous-graphe couvrant biparti]\, % 3
    Montrer que tout graphe \(G\) contient un sous-graphe biparti couvrant \(H\) avec \(m(H) \geq m(G)/2\) (on pourra ordonner les sommets de \(G\) en \(v_1,\ldots,v_n\) et construire la bipartition \(X, Y\) de \(H\) en pla\c cant \(v_i\) dans \(X\) ou \(Y\) en fonction des \(v_j\) pour \(j < i\)).
\end{td-exo}

\begin{td-sol}\,\\
    Soit \(H\) le sous graphe induit de \(G\) de degr\'e moyen maximum et de nombre de sommets minimum pour cela.

    Si \(\delta(H) \geq k + 1\), on a gagn\'e,

    Sinon, \(H\) contient \(x\) avec \(d_H(x) \leq k\).

    En particulier, on a \(\ol d(H) \geq \ol d(G) \geq 2k\).
    On a 
    \begin{equation*}
        \begin{aligned}
            &\hphantom{=} 2m(H \setminus x) \\
            &= 2m(H) - 2d_H(x) \geq 2(m(H) - k) \\
            &\geq \frac2{n(H)} \cdot \left(n(H) m(H) - n(H) k\right) \\
            &\geq \frac2{n(H)} \cdot \left(n(H) m(H) - m(H)\right) \\
            &\geq 2 \cdot \frac{m(H)}{n(H)} \cdot (n(H) - 1) \\
            &= \ol d(H) \cdot (n(H \setminus x)) \\
        \end{aligned}
    \end{equation*}
    Et donc
    \begin{equation*}
        \ol d(H \setminus x) = 2 \frac{m(H \setminus x)}{n(H \setminus x)} \geq \ol d(H).
    \end{equation*}

    Si \(\ol d(H \setminus x) > \ol d(H)\) alors, on contredit le choix de \(H\).
    Si \(\ol d(H \setminus x) = \ol d(H)\) alors, on contredit le choix de \(H\) aussi, car \(H \setminus x\) a un sommet de moins que \(H\) (et est non nul).
\end{td-sol}

\begin{td-exo}[Complexit\'e de sous-graphe]\,\\
    Quelle est la complexit\'e du probl\`eme algorithmique suivant, le graphe \(H\) \'etant fix\'e ?
    % insert algo 1 here

    Et de celui-ci ?

    % insert algo 2 here
\end{td-exo}